\chapter{Introduction}
\section{Background and motivation}
Gliomas are the most frequent malignant primary tumours of the brain. They arise from glial cells and range from slowly growing lower-grade lesions to aggressive glioblastoma, and their clinical management depends heavily on imaging. Multi-modal magnetic resonance imaging (MRI) is the standard tool: native T1-weighted (T1), contrast-enhanced T1-weighted (T1ce), T2-weighted (T2) and T2 fluid-attenuated inversion recovery (FLAIR) sequences each highlight different tissue properties, and together they reveal the oedema surrounding the tumour, the necrotic and non-enhancing core, and the actively enhancing rim \cite{menze2015,bakas2017}. Delineating these sub-regions is required for surgical planning, for defining radiotherapy target volumes and for measuring the response to treatment over time.

In routine practice the delineation is still largely manual. A radiologist or radiation oncologist traces the tumour boundary slice by slice through a volume of roughly 150 axial slices, which is time-consuming and subject to inter-observer variability; the BraTS benchmark itself was created partly because expert annotations of the same scan disagree noticeably \cite{menze2015}. Deep learning has made automated segmentation practical. The U-Net encoder--decoder \cite{ronneberger2015} and its volumetric extension, the 3D U-Net \cite{cicek2016}, are the basis of most current systems, and the self-configuring nnU-Net framework \cite{isensee2021} shows that a carefully engineered U-Net remains among the strongest approaches. Transformer-based models such as TransBTS \cite{wang2021transbts} and Swin UNETR \cite{hatamizadeh2022swin} add long-range context to these networks.

There is, however, a practical tension between context and cost. Purely two-dimensional (2D) networks are light and fast but treat every slice independently, although the true tumour changes smoothly from one slice to the next; their predictions can therefore jump between adjacent slices, producing fragmented sub-regions and inaccurate 3D boundaries. Fully three-dimensional (3D) networks capture inter-slice context directly but require large GPU memory and long training times, which limits their use where computing resources are modest. The 2.5D family sits between the two: it feeds a small stack of neighbouring slices to a 2D-style network so that some volumetric context is available at a fraction of the cost. A recent example is the Cross-Slice Transformer (CST) of Abdikenov \emph{et al.} \cite{abdikenov2026}, which learns how much each of five consecutive CT slices should contribute to the segmentation of the centre slice, and which achieved the best Dice among 2.5D models for lymph-node segmentation with fewer than 8 million parameters. That work was designed for single-channel CT and small lymph nodes; whether cross-slice attention helps multi-modal brain MRI and glioma sub-regions had not been tested. This project performs that adaptation and evaluates it rigorously.

\textbf{Relevance to the Sustainable Development Goals.} The project is aligned with United Nations Sustainable Development Goal 3 (SDG 3), \emph{Good Health and Well-being} \cite{un2015sdg}. It relates most directly to Target 3.4, which aims to reduce premature mortality from non-communicable diseases (including cancer) through prevention and treatment, and to Target 3.8, which calls for access to quality essential health-care services for all. Brain-tumour treatment quality depends on accurate and timely delineation, yet expert neuroradiology time is limited, especially in resource-constrained hospitals. The problem addressed is therefore the cost and inconsistency of glioma delineation; the project's contribution is a compact segmentation model and an evaluation showing that cross-slice context improves boundary accuracy and slice-to-slice consistency without the memory demands of a full 3D network. The intended beneficiaries are radiologists and radiation oncologists, who could review and correct an automatic draft instead of tracing from scratch, and ultimately patients, through faster and more consistent treatment planning. These clinical benefits are \emph{expected} outcomes: the work demonstrates technical performance on a public research dataset and has not been evaluated in a clinical workflow.

\section{Objectives}
The project has six objectives, each tied to an evaluation method and a deliverable.
The first objective is to design and implement a 2.5D Context-Aware Cross-Slice Transformer that fuses neighbouring-slice information for multi-modal MRI glioma segmentation; the deliverable is a working PyTorch implementation whose cross-slice module can be switched off by setting the number of slices $K=1$, which yields an otherwise identical 2D U-Net baseline.
The second objective is to build a reproducible experimental pipeline with patient-level data splits and evaluation that follows the BraTS convention of volumetric, per-patient metrics; it is evaluated by verifying that no patient appears in more than one of the training, validation and test sets and that test patients are used only once, and its deliverable is the public code repository.
The third objective is to measure the contribution of cross-slice context through a controlled comparison of $K=5$ against $K=1$ under identical training, using five-fold cross-validation over all 369 annotated BraTS 2020 patients, with the 3D Dice coefficient as the primary metric and paired Wilcoxon signed-rank tests to establish significance.
The fourth objective is to justify the training hyper-parameters systematically rather than by trial and error, through a one-factor-at-a-time design study and an automated Optuna search \cite{akiba2019}.
The fifth objective is to quantify 3D boundary quality and slice-to-slice consistency, the property that cross-slice attention is intended to improve, using the 95th-percentile Hausdorff distance (HD95) and a count of slices containing false-positive tumour.
The sixth objective is to make the pipeline run on the institutional Astra GPU cluster, so that the full-scale experiments that were impossible on a 4~GB laptop GPU could be completed; the deliverable is a set of scheduler-ready job scripts and the resulting five-fold results.

\section{Scope and constraints}
The intended users of the system are researchers and, in a future clinical setting, radiologists who would review its output. The scope covers automated segmentation of the three standard BraTS tumour regions from pre-operative multi-modal MRI, the comparison of the proposed model with a matched baseline, and the associated hyper-parameter study. Tumour grading, molecular-marker prediction, survival prediction and any clinical deployment are excluded. The work uses only the publicly available BraTS 2020 training set; the official BraTS validation and test sets do not release ground-truth masks to participants outside the challenge, so performance is estimated by internal cross-validation rather than on the official leaderboard. The approach is deliberately 2.5D: a full 3D model was outside the time budget and would also remove the very inter-slice modelling that the project sets out to study.

The project operated under significant computational constraints. Development and the preliminary design study were carried out on a laptop with an NVIDIA RTX 3050 GPU (4~GB), which forced small batches and a restricted protocol. The final experiments were run on the Astra cluster during a time-limited access window of two lab days plus one night, with at most three concurrent GPU jobs because the facility is shared. Each configuration was therefore trained once per fold with a fixed seed; variation across random seeds is not measured. Finally, the system is a research prototype and has not been validated as a medical device.

\section{Project report organization}
Chapter~2 defines the problem and the limitations of existing systems. Chapter~3 reviews twelve relevant research works and compares them in a table. Chapter~4 lists the functional, non-functional, hardware, software and data requirements. Chapter~5 presents the proposed architecture, the training and inference procedure, the hyper-parameter justification and the evolution of the system from the mid-review prototype. Chapter~6 reports the evaluation setup, criteria, results and their analysis, including limitations. Chapter~7 concludes the report and outlines future work. The appendix describes the source code and how to reproduce the experiments.

\clearpage
\chapter{Problem Definition}
\section{Problem statement}
Given a pre-operative multi-modal brain MRI examination consisting of co-registered T1, T1ce, T2 and FLAIR volumes, the task is to label every voxel as belonging to one of three nested tumour regions defined by the BraTS challenge: the \emph{whole tumour} (WT), comprising oedema, necrotic and non-enhancing core and enhancing tumour; the \emph{tumour core} (TC), comprising necrotic and non-enhancing core and enhancing tumour; and the \emph{enhancing tumour} (ET). The input to the proposed model is a stack of $K$ consecutive axial slices, each with four channels, and the output is the three-region mask of the centre slice; applying the model to every slice produces a full 3D segmentation.

The specific research question is whether attention across neighbouring slices, adapted from CT lymph-node segmentation, makes 2.5D glioma segmentation more accurate and more consistent through the volume than an otherwise identical 2D network, and whether any improvement is statistically reliable across patients. The stakeholders are clinicians who need accurate and consistent tumour volumes, patients whose treatment depends on them, and researchers who need lightweight models that can be trained without large GPU clusters. Inconsistent slice-wise segmentations have practical consequences: fragmented or jagged tumour volumes distort volume measurements, inflate boundary errors that matter for radiotherapy margins, and force clinicians to spend time correcting automatic output.

\section{Existing system and limitations}
The current clinical system is manual or semi-automatic contouring by experts. It is accurate when time allows, but it is slow, depends on scarce specialist time and shows inter-rater variability \cite{menze2015}. Among automated systems, three families are relevant, each with an evidence-based limitation.

Two-dimensional networks such as the U-Net \cite{ronneberger2015} and its attention-augmented variants segment each slice independently. They are efficient, but they ignore the anatomical continuity between neighbouring slices. In our own mid-review prototype, a 2D baseline produced on average 17.7 slices per patient that contained predicted tumour although the ground truth had none, and its boundary error (HD95) for the whole tumour was 17.9~mm, almost twice that of a cross-slice variant evaluated under the same conditions (Section~\ref{sec:evolution}).

Three-dimensional networks such as the 3D U-Net \cite{cicek2016}, nnU-Net \cite{isensee2021}, TransBTS \cite{wang2021transbts} and Swin UNETR \cite{hatamizadeh2022swin} model volumetric context directly and achieve the strongest results on BraTS, but they require large GPU memory and long training schedules. The base CST paper confirms this trade-off in a different domain: a 3D nnU-Net achieved the highest overall Dice, while the 2.5D CST matched much of its performance with a far smaller footprint \cite{abdikenov2026}. Recent glioma models listed in Chapter~3 add multi-scale aggregation, CNN--transformer hybrids or multi-task heads \cite{aslam2025,li2025deepglioseg,wen2025}, but none of them explicitly models how one slice relates to its neighbours.

Finally, the cross-slice mechanism that does model inter-slice relationships \cite{abdikenov2026} was designed for single-channel CT and very small targets, and its fusion operates on intensity statistics of raw slices. Its value for four-channel MRI and for large, heterogeneous gliomas with nested sub-regions is unknown. The gap addressed by this project is therefore the adaptation of cross-slice attention to multi-modal glioma MRI, together with a controlled, statistically tested measurement of what it contributes in accuracy, boundary quality and slice consistency.
